% Printed 518; principal PDF 525; Google PDF 529.
\par\clearpage
\noindent\makebox[0pt][l]{\fontsize{13}{16}\selectfont 518}\hfill{\small\textbf{REMARKS.}}\hfill\null\par\medskip

For the reciprocal weights, the formulas above give
\[
\begin{alignedat}{3}
 \text{for }n&=4\quad&
 \frac{1}{P}&=\frac{1}{rr}\left(\tfrac14g-\tfrac14\frac{3g+2}{2g+1}\right)&
 &=\frac{kk}{rr}\left(\tfrac14+\tfrac14\frac{kk+2}{kk(2kk+5)}\right)\\
 n&=6&
 \frac{1}{P}&=\frac{1}{rr}\left(\tfrac16g-\tfrac16\frac{5gg+4g-3}{3gg+2g-2}\right)&
 &=\frac{kk}{rr}\left(\tfrac16+\tfrac16\frac{k^4+4kk+3}{kk(3k^4+14kk+14)}\right)\\
 n&=3&
 \frac{1}{P}&=\frac{1}{rr}\left(\tfrac13g-\tfrac23\frac{2g+3}{3g+4}\right)&
 &=\frac{kk}{rr}\left(\tfrac13+\tfrac23\frac{kk+3}{kk(3kk+10)}\right)\\
 n&=5&
 \frac{1}{P}&=\frac{1}{rr}\left(\tfrac15g-\tfrac25\frac{4gg+7g}{5gg+8g-1}\right)&
 &=\frac{kk}{rr}\left(\tfrac15+\tfrac25\frac{k^4+5kk+5}{kk(5k^4+28kk+35)}\right).
\end{alignedat}
\]
\textsc{Gauss}, on the other hand, states in the preceding note: »Probable error of the result, taking the probable error of a reading $=1$:
\[
 \surd\bigl((kk+1)c^{(n)}-c^{(n-1)}\bigr)
 =k\left\delimiter"270370\relax\vphantom{\frac{1}{n}}\right.\!
 \left(\frac{1}{n}+\frac{2c^{(n)}}{nkk}\right).
\]
\indent Also, if the coefficients are denoted by $c^{(n)}$, $c^{(n-1)}$, etc., we have
\[
 (nkk+(n-2))c^{(n)}-nc^{(n-1)}=kk.\text{«}
\]

If readings are taken at the beginning and after the $1^{\text{st}}$, $(n-1)^{\text{th}}$ and $n^{\text{th}}$ observations, the value of the angle and its reciprocal weight may be derived as follows.

Let\par
\noindent
\begin{tabular}{@{}c@{\quad}p{.80\linewidth}@{}}
 $v_0$, $v_1$, $v_{n-1}$, $v_n$ & be the errors of the readings $A$, $B$, $Y$, $Z$;\\
 $\delta_1$, $\delta_{n-1}$, $\delta_n$ & be the sums of the sighting errors between the readings $A$ and $B$, $B$ and $Y$, $Y$ and $Z$;\\
 $x$ & be the value sought for the angle.
\end{tabular}\par

Let the weight of a reading error be $1$ as before, and let the weight of a single sighting error be $p$. The weights of $\delta_1$, $\delta_{n-1}$, $\delta_n$ are accordingly $\dfrac{p}{2}$, $\dfrac{p}{2(n-2)}$, $\dfrac{p}{2}$, since $\delta_1$ and $\delta_n$ each arise from $2$ settings, and $\delta_{n-1}$ from $n-2$ settings. Now put $\dfrac{p}{2}=\dfrac{1}{kk}$.

The errors satisfy the equations
\[
\makebox[\linewidth]{\hspace*{\parindent}\makebox[0pt][l]{\small 1)}\hfill
\(\displaystyle
 \begin{aligned}
  \delta_1 &= -A+B-x-v_0+v_1\\
  \delta_{n-1} &= -B+Y-(n-2)x-v_1+v_{n-1}\\
  \delta_n &= -Y+Z-x-v_{n-1}+v_n.
 \end{aligned}
\)\hfill\hspace*{\parindent}}
\]
\indent It is necessary that
\[
\makebox[\linewidth]{\hspace*{\parindent}\makebox[0pt][l]{\small 2)}\hfill
\(\displaystyle
 v_0v_0+v_1v_1+v_{n-1}v_{n-1}+v_nv_n+
 \frac{1}{kk}\left(\delta_1\delta_1+
 \frac{1}{n-2}\delta_{n-1}\delta_{n-1}+\delta_n\delta_n\right)=\min.
\)\hfill\hspace*{\parindent}}
\]
hold.

Writing for the moment
\[
 -A+B=L_1,\qquad -B+Y=L_2,\qquad -Y+Z=L_3,
\]
the normal equations, after multiplying at the same time by $kk$, are
{\small
\[
\begin{array}{@{}l@{\;}r@{\;}c@{\;}c@{\;}c@{\;}c@{\;}c@{\;}c@{}}
 & L_1+L_2+L_3 &=& nx & +v_0 & * & * & -v_n\\
 & \multicolumn{1}{l@{\;}}{L_1} &=& +x & +(kk+1)v_0 & -v_1 & * & *\\
 \text{3)} & -L_1+\dfrac{1}{n-2}L_2 &=& * & -v_0 &
  +\left(kk+1+\dfrac{1}{n-2}\right)v_1 & -\dfrac{1}{n-2}v_{n-1} & *\\[3pt]
 & -\dfrac{1}{n-2}L_2+L_3 &=& * & * &
  -\dfrac{1}{n-2}v_1 & +\left(kk+1+\dfrac{1}{n-2}\right)v_{n-1} & -v_n\\[3pt]
 & -L_3 &=& -x & * & * & -v_{n-1} & +(kk+1)v_n.
\end{array}
\]
}

% Printed 519; principal PDF 526; Google PDF 530.
\par\clearpage
\noindent\null\hfill{\small\textbf{OBSERVATIONS BY REPETITION.}}\hfill\makebox[0pt][r]{\fontsize{13}{16}\selectfont 519}\par\medskip

Adding the last $4$ equations gives
\[
 v_0+v_1+v_{n-1}+v_n=0;
\]
and the second and fifth equations therefore give
\[
 L_1-L_3=(kk+2)(v_0+v_n)=-(kk+2)(v_1+v_{n-1}).
\]

The first normal equation gives
\[
 \delta_1+\delta_{n-1}+\delta_n=0.
\]

Subtracting the fifth equation from the second and the fourth from the third gives
\[
\begin{array}{@{}l@{\;}r@{\;}r@{\;}c@{\;}c@{\;}r@{\;}l@{}}
 L_1 && +L_3 &=& 2x & +(1+kk)(v_0-v_n) & -\hphantom{\dfrac{kk(n-2)+n}{n-2}}(v_1-v_{n-1})\\
 -L_1 & +\dfrac{2}{n-2}L_2 & -L_3 &=& * & -(v_0-v_n) & +\dfrac{kk(n-2)+n}{n-2}(v_1-v_{n-1}),
\end{array}
\]
from which it follows that
\[
 (kk(n-2)+2)(L_1+L_3)+2L_2=2(kk(n-2)+n)x+
 (k^2(n-2)+2kk(n-1)+2)(v_0-v_n).
\]

Eliminating $v_0-v_n$ from the first of equations 3) by means of this equation, and at the same time substituting the values for the $L$'s, gives
\[
 x=\frac{((n-2)kk+n)(Z-A)+(n-2)(Y-B)}{n(n-2)kk+2nn-4n+4}.
\]

To obtain the reciprocal weight of $x$, replace the constants in equations 3) (because the original normal equations have been multiplied by $kk$) with $kk$, $0$, $0$, $0$, $0$, in that order. The resulting value of $x$ gives the reciprocal weight. We obtain
\[
 \frac{1}{P}=\frac{(n-2)k^4+2(n-1)kk+2}{n(n-2)kk+2nn-4n+4}.
\]

If, on the other hand, as \textsc{Gauss} did in the last formula on p.~514 and in the two following formulas on p.~515, the weight of the sighting errors belonging to one observation is set $=1$ and the weight of a reading error $=kk$, then another factor of $kk$ must be included in the denominator of the preceding expression.

For $n=12$ and $kk=\frac19$, for example, the weight of the value of the angle then becomes
\[
 P=\frac{\frac19\left\{12.10.\frac19+2.144-4.12+4\right\}}{10.\frac{1}{81}+2.11.\frac19+2}
 =6{,}26.
\]

If, however, a reading had been taken after each observation, then according to pp.~516/517, for $kk=\frac19$ we would first have
\[
 \mu_1=1,\qquad \mu_2=2{,}111,\qquad \mu_3=3{,}457,\qquad
 \mu_4=5{,}187,\qquad \mu_5=7{,}493,\qquad \mu_6=10{,}631;
\]
and, further, with $n=12$,
\[
 (nkk+n-2)\mu_m-n\mu_{m-1}=30{,}574=kkM.
\]

Since in this case the expression for the reciprocal weight must also be divided by $kk$, and since in addition $r=1$, this gives
\[
 \frac{1}{P}=\frac{1}{n}+\frac{2\mu_m}{nkkM}
\]
or

% Printed 520; principal PDF 527; Google PDF 531.
\par\clearpage
\noindent\makebox[0pt][l]{\fontsize{13}{16}\selectfont 520}\hfill{\small\textbf{REMARKS.\quad OBSERVATIONS BY REPETITION.}}\hfill\null\par\medskip
\[
 P=\frac{12kkM}{kkM+2\mu_6}=\frac{12.30{,}574}{51{,}836}=7{,}08.
\]
\indent Compare the end of the letter to \textsc{Bessel}, p.~516.

If readings are taken only at the beginning and after the $n^{\text{th}}$ observation, then, either directly as before or by setting $n=2m+1=1$ in the general formula on p.~517, we find for the reciprocal weight of $x=\dfrac{Z-A}{n}$:
\[
 \frac{1}{P}=\frac{nkk+2}{nn}.
\]

Here the weight of a reading error is set $=1$, and the weight of the total pointing error is set $=\dfrac{1}{nkk}$.

\textsc{Svanberg}'s method (Exposition des opérations faites en Laponie, pour la détermination d'un arc du méridien, en 1801, 1802 et 1803, Stockholm 1805, p.~29/33) consists in forming, to derive the value $x$ of the angle from $(n+1)$ readings $l_0$, $l_1$, \ldots, $l_n$, all possible differences $l_i-l_k=(i-k)x$, $i>k$. Their sum gives
\[
 \frac{n(n+1)(n+2)}{1.2.3}x=\sum_{i=0}^{n}(2i-n)l_i.
\]

As \textsc{Gauss} notes in the preceding letter to \textsc{Bessel}, the same result is obtained by setting the pointing errors, $\delta$, equal to zero. In fact, this value of $x$ also follows from the error equations
\[
 l_0+v_0=u,\quad l_1+v_1=u+x,\quad l_2+v_2=u+2x,\quad \ldots,\quad l_n+v_n=u+nx,
\]
where $u$ denotes the distance from the origin of the graduations to the first reading position.\par
\hfill\textsc{Krüger.}\hspace*{\parindent}\par

\vspace{2.5cm}
\begin{center}\rule{4cm}{.4pt}\end{center}

% Unfoliated [521]; principal PDF 528; Google PDF 532.
\par\clearpage
\vspace*{.11\textheight}
\begin{center}
 {\Large REMARKS ON THE NINTH VOLUME.}\par
 \vspace{1em}
 \rule{1.5cm}{.4pt}
\end{center}
\vspace{1em}

The present ninth volume of \textsc{Gauss}'s works was edited by Mr. L.~\textsc{Krüger} in Potsdam, with the assistance of Mr. A.~\textsc{Börsch} in examining the material. It contains both reprints of the geodetic writings published by \textsc{Gauss} for which there was no room in Volume IV and an edition of the posthumous geodetic papers.

The geodetic notes left by \textsc{Gauss} concern almost exclusively the \textit{Hanoverian arc measurement} and questions prompted by it and by its predecessor, the Danish arc measurement. Mr. L.~\textsc{Krüger} undertook the laborious task of arranging into a coherent whole the contents of the disconnected leaves among the posthumous papers on which \textsc{Gauss} developed the \textit{projection method for the Hanoverian arc measurement}, as well as the records concerning the \textit{triangulation} itself. A collection of all the important publications, official reports, and passages from letters on the subject gives a detailed account of the \textit{progress of the work on the arc measurement and the subsequent land survey}. For details, see the editor's remarks appended to the various notes. An overview of the principal parts of this volume may also be found in the »Bericht über den Stand der Herausgabe von \textsc{Gauss}' Werken« presented by the undersigned in April 1902 (Nachrichten der K.~Gesellschaft der Wissenschaften zu Göttingen. Geschäftliche Mittheilungen 1902, Heft 1).

As with Volume VIII, the general editorship was in the hands of Mr. M.~\textsc{Brendel}. The excerpts from the posthumous papers are printed as closely as possible to the originals. As in the earlier volumes, the spellings that \textsc{Gauss} appeared to have used most frequently were chosen and retained throughout the volume. As before, the editor's insertions in the original text are marked by square brackets [] or mentioned specifically in the remarks, and reprints of notes and passages from letters that do not originate with \textsc{Gauss} are placed in braces \{\}.

The passages from letters were printed from the originals wherever these survive. The source of the notes printed from the posthumous papers is given, along with other editorial information, in the remarks to the individual sections.\par
\hfill F.~\textsc{Klein.}\hspace*{\parindent}\par

\begin{center}\rule{4cm}{.4pt}\end{center}
\noindent\makebox[0pt][l]{\small IX.}\makebox[.86\linewidth][r]{66}\par

% Unfoliated [522]; principal PDF 529; Google PDF 533.
\par\clearpage
\begingroup
\small
\newcount\IXContentsDot
\newsavebox{\IXContentsLabel}
\newlength{\IXContentsAvailable}
\newlength{\IXContentsPeriod}
\settowidth{\IXContentsPeriod}{.}
\newcommand{\IXContentsFit}[2]{%
  \sbox{\IXContentsLabel}{#2}%
  \ifdim\wd\IXContentsLabel>#1\relax
    \resizebox{#1}{!}{\usebox{\IXContentsLabel}}%
  \else
    \usebox{\IXContentsLabel}%
  \fi
}
\newcommand{\IXContentsEntry}[5]{%
  \par\noindent\hbox to \linewidth{%
    \setlength{\IXContentsAvailable}{\dimexpr\linewidth-#1-4.15em\relax}%
    \ifnum#3>0\relax
      \addtolength{\IXContentsAvailable}{\dimexpr-.9em-\IXContentsPeriod*#3\relax}%
    \fi
    \hspace*{#1}\strut\IXContentsFit{\IXContentsAvailable}{#2}%
    \ifnum#3>0\relax
      \hskip.35em
      \IXContentsDot=0\relax
      \loop\ifnum\IXContentsDot<#3\relax
        \hfil.\advance\IXContentsDot by 1\relax
      \repeat
      \hskip.55em
    \fi
    \hfil\makebox[2.35em][c]{#4}\makebox[1.8em][r]{#5}%
  }\par
}
\newcommand{\IXContentsText}[2]{%
  \par\noindent\hspace*{#1}\strut\IXContentsFit{\dimexpr\linewidth-#1\relax}{#2}\par
}
\newcommand{\IXContentsHeading}[1]{%
  \par\medskip\noindent\makebox[\linewidth][c]{\IXContentsFit{\linewidth}{#1}}\par\smallskip
}

\vspace*{.11\textheight}
\begin{center}
 {\Large\textls[200]{CONTENTS.}}\par
 \vspace{1em}
 {\large GAUSS'S WORKS, VOLUME IX.\\
 GEODESY.\quad CONTINUATION OF VOLUME IV.}\par
 \vspace{1em}
 \rule{1.5cm}{.4pt}
\end{center}
\vspace{1em}
\begin{center}
 DETERMINATION OF THE DIFFERENCE IN LATITUDE BETWEEN THE OBSERVATORIES\\
 OF GÖTTINGEN AND ALTONA BY OBSERVATIONS\\
 WITH RAMSDEN'S ZENITH SECTOR.
\end{center}
\IXContentsEntry{1.8em}{Introduction}{29}{Page}{5}
\IXContentsEntry{1.8em}{The observed stars}{24}{—}{8}
\IXContentsEntry{1.8em}{The observations}{26}{—}{10}
\IXContentsText{1.8em}{Results}
\IXContentsText{3.6em}{Simplest combination of the observations to determine the difference in latitude,}
\IXContentsEntry{5.4em}{Art.~1}{28}{—}{29}
\IXContentsEntry{3.6em}{Accuracy of the observations, Art.~2}{17}{—}{30}
\IXContentsEntry{3.6em}{Collimation error, Art.~3}{23}{—}{31}
\IXContentsEntry{3.6em}{Absolutely most advantageous combination of the observations, Art.~4—7}{7}{—}{32}
\IXContentsEntry{3.6em}{Allowance for irregular errors of graduation, Art.~8—11}{8}{—}{36}
\IXContentsEntry{3.6em}{Locations of the observing sites, Art.~12}{18}{—}{40}
\IXContentsText{3.6em}{Determination of the absolute altitude of the pole at Göttingen Observatory from observations of the}
\IXContentsEntry{5.4em}{Pole Star with \textsc{Reichenbach}'s meridian circle, Art.~13—17}{7}{—}{40}
\IXContentsEntry{3.6em}{Final result of the Hanoverian arc measurement, Art.~18—20}{11}{—}{47}
\IXContentsText{3.6em}{Comparison of the declinations of the observed zenith stars with \textsc{Bradley}'s and}
\IXContentsEntry{5.4em}{\textsc{Piazzi}'s determinations, Art.~21}{19}{—}{50}
\IXContentsEntry{1.8em}{Determination of the latitude of Seeberg Observatory}{18}{—}{52}
\IXContentsEntry{1.8em}{Supplement to Art.~20, p.~48}{24}{—}{56}
\IXContentsText{0pt}{\textit{Notice.}}
\IXContentsText{1.8em}{Determination of the difference in latitude between the observatories of Göttingen and Altona}
\IXContentsEntry{3.6em}{by observations with \textsc{Ramsden}'s zenith sector}{12}{—}{59}
\IXContentsEntry{1.8em}{Remarks}{28}{—}{63}

% Printed 523; principal PDF 530; Google PDF 534.
\par\clearpage
\noindent\null\hfill{\small\textbf{CONTENTS.}}\hfill\makebox[0pt][r]{\fontsize{13}{16}\selectfont 523}\par\medskip

\IXContentsHeading{EARTH ELLIPSOID AND GEODESIC.}
\IXContentsText{0pt}{\textit{Posthumous papers.}}
\IXContentsEntry{1.8em}{The Earth ellipsoid}{27}{Page}{67}
\IXContentsEntry{1.8em}{Equation of the vertical plane of the ellipsoid of revolution}{15}{—}{70}
\IXContentsEntry{1.8em}{Equation of the ellipsoid of revolution relative to a tangent plane}{6}{—}{70}
\IXContentsEntry{1.8em}{Remarks}{28}{—}{71}
\IXContentsEntry{1.8em}{Foundation of my theory of the geodesic}{15}{—}{72}
\IXContentsEntry{1.8em}{Shortest line on the spheroid}{21}{—}{74}
\IXContentsEntry{1.8em}{Geodesic}{26}{—}{78}
\IXContentsEntry{1.8em}{Geodetic transfer of latitude, longitude and azimuth}{12}{—}{80}
\IXContentsEntry{1.8em}{Geodetic transfer on the sphere}{19}{—}{88}
\IXContentsText{1.8em}{Calculation of the linear length of the geodesic and its azimuths from the}
\IXContentsEntry{3.6em}{geographic coordinates}{23}{—}{89}
\IXContentsEntry{1.8em}{Perfectly exact formulas for a triangle on the elliptic spheroid}{6}{—}{92}
\IXContentsText{1.8em}{Transfer of geographic position by means of the chord and the azimuth of the}
\IXContentsEntry{3.6em}{vertical section}{28}{—}{93}
\IXContentsEntry{1.8em}{The difference between the geodetic and the observed azimuth}{6}{—}{94}
\IXContentsEntry{1.8em}{Reduction of the astronomical azimuth to the geodetic azimuth}{11}{—}{95}
\IXContentsEntry{1.8em}{Remarks}{28}{—}{96}
\IXContentsText{0pt}{\textit{Correspondence.}}
\IXContentsEntry{1.8em}{Change in the altitude of the pole with height: \textsc{Gauss} to J.~J. \textsc{Baeyer} 1853 June 22}{6}{—}{99}
\IXContentsEntry{1.8em}{Remarks}{28}{—}{102}
\IXContentsText{0pt}{\textit{Posthumous papers.}}
\IXContentsEntry{1.8em}{Reduction of spherical triangle angles $A$, $B$, $C$ to chord angles $\mathfrak{A}$, $\mathfrak{B}$, $\mathfrak{C}$}{3}{—}{103}
\IXContentsText{1.8em}{Condition for 3 points on the surface of a sphere to lie}
\IXContentsEntry{3.6em}{on a great circle}{29}{—}{103}
\IXContentsEntry{1.8em}{Remarks}{28}{—}{104}

\IXContentsHeading{CONFORMAL DOUBLE PROJECTION OF THE SPHEROID ONTO THE SPHERE AND THE PLANE.}
\IXContentsText{0pt}{\textit{Posthumous papers.}}
\IXContentsEntry{1.8em}{Mapping the elliptic spheroid onto the sphere}{15}{—}{107}
\IXContentsEntry{1.8em}{Remarks}{28}{—}{115}
\IXContentsEntry{1.8em}{Stereographic projection of the sphere onto the plane}{14}{—}{117}
\IXContentsEntry{1.8em}{Remarks}{28}{—}{121}
\IXContentsEntry{1.8em}{Mapping the sphere onto the plane by \textsc{Mercator}'s projection}{8}{—}{123}
\IXContentsEntry{1.8em}{Remarks}{28}{—}{132}
\medskip
\IXContentsEntry{1.8em}{Stereographic representation of the spheroid in the plane}{12}{—}{133}
\IXContentsEntry{1.8em}{Remarks}{28}{—}{134}

\vspace{1.7cm}
\noindent\makebox[.86\linewidth][r]{66\textsuperscript{*}}\par

% Printed 524; principal PDF 531; Google PDF 535.
\par\clearpage
\noindent\makebox[0pt][l]{\fontsize{13}{16}\selectfont 524}\hfill{\small\textbf{CONTENTS.}}\hfill\null\par\medskip

\IXContentsHeading{CONFORMAL MAPPING OF THE SPHEROID ONTO THE CONICAL SURFACE.}
\IXContentsText{0pt}{\textit{Posthumous papers.}}
\IXContentsEntry{1.8em}{On the second representation of the spheroid, referred to a parallel of latitude}{6}{Page}{137}
\IXContentsEntry{1.8em}{Remarks}{28}{—}{140}

\IXContentsHeading{CONFORMAL MAPPING OF THE SPHEROID IN THE PLANE}
\noindent\makebox[\linewidth][c]{(PROJECTION METHOD OF THE HANOVERIAN LAND SURVEY).}\par\smallskip
\IXContentsText{0pt}{\textit{Posthumous papers.}}
\IXContentsEntry{1.8em}{Calculation of geographic latitude and longitude from plane rectangular coordinates}{0}{—}{143}
\IXContentsEntry{1.8em}{Calculation of meridian convergence from plane rectangular coordinates}{4}{—}{146}
\IXContentsEntry{1.8em}{Formulas for the numerical calculation of longitude, latitude and meridian convergence}{4}{—}{148}
\IXContentsEntry{1.8em}{Calculation of the scale factor $n$}{17}{—}{152}
\IXContentsEntry{1.8em}{Relations between $x$, $y$ and $\xi$, $\lambda$}{20}{—}{155}
\IXContentsEntry{1.8em}{Calculation of plane rectangular coordinates from geographic latitude and longitude}{0}{—}{156}
\IXContentsEntry{1.8em}{Calculation of meridian convergence from geographic coordinates}{7}{—}{158}
\IXContentsEntry{1.8em}{Reduction of the azimuth on the spheroid to the azimuth in plano}{6}{—}{159}
\IXContentsText{1.8em}{The difference between the projection of the geodesic and the straight line}
\IXContentsText{3.6em}{joining its endpoints in the conformal mapping of a curved surface}
\IXContentsEntry{3.6em}{in the plane}{29}{—}{162}
\IXContentsEntry{1.8em}{On the transformation of coordinates}{20}{—}{168}
\IXContentsEntry{1.8em}{Series relating $\varphi$, $\psi$ and $\omega$}{22}{—}{171}
\IXContentsEntry{1.8em}{On calculating $\log\cos\varphi$}{22}{—}{180}
\IXContentsEntry{1.8em}{Calculation of $\log(1-ee\sin\varphi^2)$, $\log\frac{(1-ee\sin\varphi^2)^2}{aa(1-ee)}$ and $\frac{a\cos\varphi}{\surd(1-ee\sin\varphi^2)}$}{7}{—}{181}
\IXContentsEntry{1.8em}{Numerical values of the coefficients in the series relating $\varphi$, $\psi$ and $\omega$}{6}{—}{182}
\IXContentsText{1.8em}{Calculation of plane rectangular coordinates from geographic coordinates by}
\IXContentsEntry{3.6em}{means of the series relating $\varphi$, $\psi$ and $\omega$}{17}{—}{185}
\IXContentsEntry{1.8em}{Calculation of longitude and latitude from plane coordinates}{10}{—}{191}
\IXContentsEntry{1.8em}{Representation of the spheroid's surface in the plane}{11}{—}{193}
\IXContentsEntry{1.8em}{Remarks}{28}{—}{195}
\IXContentsText{0pt}{\textit{Correspondence.}}
\IXContentsText{1.8em}{On the formulas for the Hanoverian land survey:}
\IXContentsEntry{.36\linewidth}{\textsc{Gauss} to \textsc{Schumacher} 1830 April 18}{7}{—}{205}
\IXContentsEntry{.36\linewidth}{\textsc{Gauss} to \textsc{Schumacher} 1830 April 30}{7}{—}{212}
\IXContentsEntry{.36\linewidth}{\textsc{Gauss} to \textsc{Schumacher} 1831 May 17}{7}{—}{213}
\IXContentsEntry{.36\linewidth}{\textsc{Gauss} to \textsc{Schumacher} 1831 June 25}{7}{—}{215}
\IXContentsEntry{.36\linewidth}{\textsc{Gauss} to \textsc{Schumacher} 1838 Dec.~9}{7}{—}{217}
\IXContentsEntry{1.8em}{Remarks}{28}{—}{218}

\IXContentsHeading{TRIGONOMETRIC POINT DETERMINATION.}
\IXContentsText{0pt}{\textit{Posthumous papers.}}
\IXContentsText{1.8em}{Final result for the position of a point in a plane fixed}
\IXContentsEntry{3.6em}{by intersection from three known points}{27}{—}{221}

% Printed 525; principal PDF 532; Google PDF 536.
\par\clearpage
\noindent\null\hfill{\small\textbf{CONTENTS.}}\hfill\makebox[0pt][r]{\fontsize{13}{16}\selectfont 525}\par\medskip

\IXContentsText{1.8em}{Determination of the position of a point $P^\circ$ from the positions of three others, $P$, $P'$, $P''$, at which it is}
\IXContentsEntry{3.6em}{observed}{27}{Page}{223}
\IXContentsEntry{1.8em}{Adjustment of three intersections}{23}{—}{224}
\IXContentsEntry{1.8em}{On the adjustment of three intersections}{21}{—}{226}
\IXContentsText{1.8em}{Determination of a subsidiary point (Scheessel) from observations at principal triangulation points}
\IXContentsEntry{3.6em}{(Litberg, Wilsede, Bottel, Bullerberg and Brüttendorf)}{12}{—}{228}
\IXContentsText{0pt}{\textit{Treatise.}}
\IXContentsEntry{1.8em}{Application of probability theory to a problem of practical geometry}{1}{—}{231}
\IXContentsText{0pt}{\textit{Posthumous papers.}}
\IXContentsText{1.8em}{Determination of the position of point $X$ by observing the angles $\alpha$ and $\beta$ between}
\IXContentsEntry{3.6em}{3 given points $A$, $B$, $C$}{21}{—}{238}
\IXContentsEntry{1.8em}{Orientation of the plane table}{23}{—}{239}
\IXContentsEntry{1.8em}{Problem of practical geometry}{20}{—}{239}
\IXContentsEntry{1.8em}{Remarks}{28}{—}{240}

\IXContentsHeading{ADJUSTMENT OF SIMPLE FIGURES.}
\IXContentsText{0pt}{\textit{Posthumous papers.}}
\IXContentsEntry{1.8em}{Adjustment of a quadrilateral}{23}{—}{245}
\IXContentsEntry{1.8em}{Equation relating the sides and diagonals of a quadrilateral}{10}{—}{248}
\IXContentsText{0pt}{\textit{Correspondence.}}
\IXContentsText{1.8em}{On the choice of the condition equation from the side ratios:}
\IXContentsEntry{.36\linewidth}{\textsc{Gauss} to \textsc{Gerling} 1824 Feb.~11}{8}{—}{249}
\IXContentsEntry{.36\linewidth}{\textsc{Gauss} to \textsc{Gerling} 1840 Jan.~19}{9}{—}{250}
\IXContentsText{0pt}{\textit{Posthumous papers.}}
\IXContentsEntry{1.8em}{On the adjustment of the angles in the quadrilateral}{18}{—}{254}
\IXContentsEntry{1.8em}{Quadrilateral between 4 points 1.~2.~3.~4}{19}{—}{257}
\IXContentsEntry{1.8em}{Adjustment of a polygon}{23}{—}{257}
\IXContentsEntry{1.8em}{Weight of height determinations}{21}{—}{258}
\IXContentsEntry{1.8em}{Remarks}{28}{—}{259}

\IXContentsHeading{STATION ADJUSTMENTS.}
\IXContentsText{0pt}{\textit{Posthumous papers.}}
\IXContentsText{1.8em}{Station adjustment for Zeven from all the measurements of 1824 and 1825 (excluding those}
\IXContentsEntry{3.6em}{of 4 and 5 August)}{23}{—}{263}
\IXContentsEntry{1.8em}{Station adjustment for Brillit}{22}{—}{265}
\IXContentsEntry{1.8em}{Wilsede from all the measurements of 1822 and 1824}{12}{—}{267}
\IXContentsEntry{1.8em}{Adjustment of the angles measured on the Windberg}{12}{—}{271}
\IXContentsEntry{1.8em}{Observations on Breithorn 1822}{20}{—}{274}
\IXContentsText{0pt}{\textit{Correspondence.}}
\IXContentsEntry{1.8em}{\rlap{On station adjustments:}\hspace{.32\linewidth}\textsc{Gauss} to \textsc{Gerling} 1823 Dec.~26}{10}{—}{278}
\IXContentsEntry{.36\linewidth}{\textsc{Gauss} to \textsc{Schumacher} 1827 Dec.~22}{8}{—}{281}
\IXContentsEntry{.36\linewidth}{\textsc{Gauss} to \textsc{Schumacher} 1828 Jan.~7}{9}{—}{285}
\IXContentsEntry{1.8em}{Remarks}{28}{—}{287}

% Printed 526; principal PDF 533; Google PDF 537.
\par\clearpage
\noindent\makebox[0pt][l]{\fontsize{13}{16}\selectfont 526}\hfill{\small\textbf{CONTENTS.}}\hfill\null\par\medskip

\IXContentsHeading{ON NETWORK ADJUSTMENT.}
\IXContentsText{0pt}{\textit{Posthumous papers.}}
\IXContentsEntry{1.8em}{Number of condition equations in a system of triangles}{11}{Page}{297}
\IXContentsEntry{1.8em}{The 33 principal triangulation points and their connections}{14}{—}{297}
\IXContentsEntry{1.8em}{Compilation of the observed triangles and their discrepancies}{8}{—}{300}
\IXContentsEntry{1.8em}{Normal equations corresponding to the angle equations}{11}{—}{302}
\IXContentsEntry{1.8em}{Condition equations of the second kind}{19}{—}{304}
\IXContentsEntry{1.8em}{Normal equations corresponding to the side equations}{12}{—}{311}
\IXContentsEntry{1.8em}{The corrections}{26}{—}{312}
\IXContentsEntry{1.8em}{Adjusted values}{26}{—}{314}
\IXContentsEntry{1.8em}{Azimuths of the sides of the spheroidal and plane triangle systems}{5}{—}{317}
\IXContentsText{0pt}{\textit{Correspondence.}}
\IXContentsEntry{1.8em}{\textsc{Gauss} to \textsc{Olbers} 1823 Nov.~2}{22}{—}{319}
\IXContentsEntry{1.8em}{\textsc{Gauss} to \textsc{Olbers} 1824 July 6}{22}{—}{320}
\IXContentsEntry{1.8em}{\textsc{Gauss} to \textsc{Olbers} 1826 May 14}{22}{—}{320}
\IXContentsEntry{1.8em}{\textsc{Gauss} to \textsc{Gerling} 1837 May 2}{22}{—}{323}
\IXContentsEntry{1.8em}{\textsc{Gauss} to \textsc{Gerling} 1838 June 5}{21}{—}{323}
\IXContentsEntry{1.8em}{\textsc{Gauss} to \textsc{Gerling} 1838 Nov.~14}{21}{—}{326}
\IXContentsEntry{1.8em}{Remarks}{28}{—}{327}

\IXContentsHeading{RING OF TRIANGLES AROUND OLDENBURG.}
\IXContentsText{0pt}{\textit{Posthumous papers.}}
\IXContentsEntry{1.8em}{On the adjustment of the ring of triangles surrounding the territory of Oldenburg}{8}{—}{331}
\IXContentsEntry{1.8em}{Remarks}{28}{—}{340}

\IXContentsHeading{ON THE HANOVER TRIANGULATION.}
\IXContentsText{0pt}{\textit{Correspondence.}}
\IXContentsText{1.8em}{\textsc{Gauss} to \textsc{Schumacher}: 1816 July 5, 1818 Sept.~10, 1820 May 20, 1821 Mar.~4, 1822 Sept.~18,}
\IXContentsEntry{3.6em}{1825 January, 1825 June 20, 1829 Jan.~14}{17}{—}{345}
\IXContentsText{1.8em}{\textsc{Gauss} to \textsc{Bessel}: 1821 Dec.~26, 1822 Nov.~15, 1823 Nov.~5, 1824 Nov.~20, 1826 Mar.~12,}
\IXContentsEntry{3.6em}{1826 Nov.~20, 1827 April 1, 1830 April 9}{16}{—}{349}
\IXContentsEntry{1.8em}{\textsc{Gauss} to \textsc{Bohnenberger}: 1823 Nov.~16}{18}{—}{364}
\IXContentsText{1.8em}{\textsc{Gauss} to \textsc{Olbers}: 1821 Jan.~13, 1822 April 18, 1824 July 4, 1824 July 8, 1825 Feb.~19,}
\IXContentsText{3.6em}{1825 Feb.~25, 1825 Oct.~9, 1826 April 2, 1827 Jan.~14, 1827 Mar.~1, 1830 June 14,}
\IXContentsEntry{3.6em}{1837 Sept.~2}{27}{—}{367}
\IXContentsText{1.8em}{\textsc{Gauss} to \textsc{Gerling}: 1821 Oct.~5, 1822 Feb.~21, 1822 Nov.~7, 1823 July 27, 1823 Aug.~11,}
\IXContentsEntry{3.6em}{1823 Sept.~1, 1823 Sept.~5, 1823 Oct.~3, 1827 July 19, 1838 Sept.~12, 1838 Nov.~14}{2}{—}{380}
\IXContentsEntry{1.8em}{Remarks}{28}{—}{395}
\IXContentsText{0pt}{\textit{Publications.}}
\IXContentsEntry{1.8em}{Extract from a letter by Court Councillor and Knight \textsc{Gauss}}{10}{—}{396}
\IXContentsEntry{1.8em}{Extract. from a letter by Court Councillor \textsc{Gauss}}{13}{—}{397}
\IXContentsEntry{1.8em}{Remarks}{28}{—}{400}
\IXContentsText{0pt}{\textit{Posthumous papers.}}
\IXContentsText{1.8em}{Plan and opening of the work on the trigonometric surveys in Hanover.}
\IXContentsEntry{3.6em}{Plan of the work}{25}{—}{401}

% Printed 527; principal PDF 534; Google PDF 538.
\par\clearpage
\noindent\null\hfill{\small\textbf{CONTENTS.}}\hfill\makebox[0pt][r]{\fontsize{13}{16}\selectfont 527}\par\medskip

\IXContentsEntry{3.6em}{Introduction}{28}{Page}{402}
\IXContentsEntry{3.6em}{First Section.\quad General arrangement of the surveys}{9}{—}{405}
\IXContentsText{1.8em}{Extracts from reports on the triangulation to the Hanoverian Cabinet Ministry.}
\IXContentsEntry{3.6em}{From a report of 7 January 1822 on the work in 1821}{5}{—}{406}
\IXContentsText{3.6em}{From a report of 31 January 1823 on the arc measurement work in}
\IXContentsEntry{5.4em}{1822}{28}{—}{409}
\IXContentsText{3.6em}{From the report of 16 February 1825 on the trigonometrical work}
\IXContentsEntry{5.4em}{carried out in 1824}{23}{—}{412}
\IXContentsText{3.6em}{From the report of 21 November 1827 on the further extension of the}
\IXContentsEntry{5.4em}{arc measurement work}{24}{—}{413}
\IXContentsEntry{3.6em}{From the report of 26 June 1828 on the continuation of the arc measurement work}{0}{—}{416}
\IXContentsText{3.6em}{Historical report on the surveys partly carried out and partly directed}
\IXContentsEntry{5.4em}{by Court Councillor \textsc{Gauss} in the Kingdom of Hanover}{14}{—}{418}
\IXContentsText{3.6em}{From the report of 5 July 1840 on the trigonometrical surveys in}
\IXContentsEntry{5.4em}{1839}{28}{—}{425}
\IXContentsText{3.6em}{From a report of December 1844 on the trigonometrical measurements}
\IXContentsEntry{5.4em}{carried out in 1844}{22}{—}{425}
\IXContentsEntry{1.8em}{Geographic coordinates of the principal triangulation points of the Hanoverian measurements}{3}{—}{427}
\IXContentsEntry{1.8em}{Remarks}{28}{—}{429}

\bigskip
\IXContentsHeading{HEIGHT MEASUREMENTS.}
\IXContentsText{0pt}{\textit{Publication and correspondence.}}
\IXContentsText{1.8em}{The refraction coefficient from the height measurements in the Hanoverian arc measurement.}
\IXContentsText{3.6em}{Observed and computed triangulation in the Hanover, Braunschweig, and}
\IXContentsEntry{5.4em}{Lüneburg territories}{24}{—}{437}
\IXContentsText{3.6em}{The coefficient of refraction derived from the height measurements of 1824: \textsc{Gauss} to}
\IXContentsEntry{5.4em}{\textsc{Schumacher} 1824 Nov.~28}{20}{—}{440}
\IXContentsText{3.6em}{The coefficient of refraction from the height measurements of the arc measurement and}
\IXContentsEntry{5.4em}{its continuation to Jever: \textsc{Gauss} to \textsc{Schumacher} 1846 Dec.~27}{8}{—}{443}
\IXContentsEntry{3.6em}{Remarks}{26}{—}{444}
\IXContentsText{0pt}{\textit{Posthumous papers.}}
\IXContentsEntry{1.8em}{Terrestrial refraction}{25}{—}{445}
\IXContentsEntry{1.8em}{Formula for calculating heights}{22}{—}{448}
\IXContentsEntry{1.8em}{Adjustment of the heights of the principal triangulation points}{15}{—}{449}
\IXContentsEntry{1.8em}{Remarks}{28}{—}{454}
\IXContentsText{0pt}{\textit{Publication.}}
\IXContentsEntry{1.8em}{Table for barometric height measurement}{19}{—}{456}

\bigskip
\IXContentsHeading{HELIOTROPE.}
\IXContentsText{0pt}{\textit{Publications.}}
\IXContentsEntry{1.8em}{On the Heliotrope}{26}{—}{461}
\IXContentsEntry{1.8em}{Invention of a Heliotrope}{23}{—}{466}
\endgroup
